\section*{Technical Validation}

%%% NOTES. (journal instructions)
% This section presents any experiments or analyses that are needed to support the technical quality of the dataset. This section may be supported by figures and tables, as needed. This is a required section; authors must present information justifying the reliability of their data.

This section motivates and details the parameter choices for the source retrieval and text alignment components of the text reuse detection pipeline. For both steps, we strive for maximum accuracy given the constraints for scalability and computational efficiency imposed by the amount of data to be processed. 


\subsection*{Source Retrieval}

Objective of the source retrieval step is to prune the search space of document pairs by reducing the number of pairs to be compared in subsequent (computationally expensive) steps. The optimization criterion is recall: Ideally, no document pair containing text reuse should be overlooked, i.e., the false negative rate should be minimized.

To allow for a fine-grained detection of local similarity, the parameters of the source retrieval step are set to $n=50$ and $m=10$, i.e., documents are split into chunks of 50~words, with 10~MinHash values computed per chunk. Consequently, the source retrieval is able to identify pairs of documents that share as little as a 9-word overlap between any two chunks. Note that these words do not need to be consecutive since chunks are represented in a bag-of-words model. Since the subsequent alignment step operates at the 8-gram level (=~eight consecutive units, see next section), the filtering within the source retrieval step eliminates only pairs for which the subsequent alignment step is guaranteed to not find any matches. Overall, the source retrieval step yields $3.305\times 10^{12}$ total unique document pairs for further analysis. Given the initial document count of 4,656,302, this represents a pruning of the search space by over~84\% compared to an exhaustive pairwise comparison, rendering the source retrieval step highly effective in improving overall efficiency.


\subsection*{Text Alignment}

The text alignment step identifies the exact location and extent of reuse between two documents. Unlike source retrieval, text alignment is a precision-oriented task: Special emphasis is put on precision over recall for the parameter choice to minimize the number of false positives; a low noise ratio is paramount for meaningful future analyses of the dataset. To ensure the effectiveness of the text alignment, its parameters are chosen by performing a grid-search, using precision, recall, and $F_{0.5}$ as effectiveness scores. We employ the PAN-13 Text Alignment Corpus for evaluation, which has been previously published as a benchmark dataset for text alignment by the PAN Shared Task on Plagiarism Detection \cite{potthast:2013}. It contains $10,000$~pairs of texts that are subject to different so-called obfuscation strategies that simulate more difficult cases of plagiarism, where the authors tried to hide text-reuse by paraphrasing the copied text to some extent. 

The applied obfuscation strategies include \emph{random obfuscation} (shuffling, adding, deleting, and replacing words or short phrases at random), \emph{cyclic translation obfuscation} (a text is translated into another language, and back to the original language; possibly with more languages in-between), \emph{summary obfuscation} (human-generated summaries of source texts), non-obfuscated plagiarism, and pairs without any text reuse between source and target. Each of these strategies is equally represented in the corpus with $2,000$ pairs. The grid search identifies an \emph{n}-gram size of $n = 8$, an \emph{n}-gram overlap of $k = 7$, and an extension range of $\Delta = 250$ as optimal. Under these conditions, our implementation of the text alignment step achieves a precision of~0.93 at a recall of~0.46, and thus an~$F_{0.5}$ score of~0.77. This renders our approach as highly competitive when compared to other approaches evaluated as part of the PAN Shared Tasks, placing it among the best for precision and $F_{0.5}$, while being the only approach adhering to the scalability requirements imposed by the scale of our analysis. A full overview of the attained scores for comparison with competing systems is given in Table~\ref{tab:pan13}.

When evaluating separately per obfuscation strategy (Table~\ref{tab:evaluation}), precision is very high throughout, exceeding~0.88 in all cases. This fact places our system within~0.08 for the best-performing, yet computationally complex approach for each obfuscation strategy. Yet, a drop in recall can be observed for heavily obfuscated text, down to~0.1 for summary obfuscation. This effect is expected and also noticeable among the other approaches presented at~PAN. Moreover, our focus is not on plagiarism but on establishing a general baseline for text reuse in science, where more literal reuse (e.g., citations, idiomatic reuse, template writing, etc.) is expected to be the norm. The heavily obfuscated test sets studied at~PAN were dedicated to study extreme cases of plagiarism, where an author expends much effort to hide the fact via severe forms of paraphrasing. The investment of such an effort, however, goes against the time savings that can be expected from reusing text, so that it can be presumed that the vast amount of paraphrased text reuse will only make smaller changes to a text instead of changing every last \emph{n}-gram. This observation is partially corroborated when reviewing the many cases of academic plagiarism found in dissertation theses throughout the recent years \cite{vroniplag:2021}, where heavy paraphrasing is hardly ever observed. We therefore deemed the investment of an exponentially higher computation cost into retrieving such cases to be uneconomical.

Nevertheless, two measures to tackle this recall issue have been proposed by PAN~participants:
\Ni
customized approaches for each of the different obfuscation strategies, employing heuristics to detect which kind a document pair is exhibiting, and
\Nii
ensemble methods encompassing different seeding and extension strategies, combined into a single result.
However, the first is not applicable to our situation, as there is no ground truth data available to fine-tune such a classification to scientific writing. The second comes at a very high runtime and algorithmic complexity.
This is reflected in Table~\ref{tab:pan13} as well: for each approach, the (asymptotic) algorithmic complexity is noted, for two given sequences of length $n$ and $m$. Only six of the approaches presented at PAN perform in sub-quadratic time, a necessary requirement for large-scale detection. Furthermore, the data specificity for each approach is listed: it denotes how much fine-tuning to the PAN data has taken place, for example by crafting specialized corpus-dependent features, using ensemble methods, trained classifiers, or specific approaches for each of the different obfuscation strategies, all of which reduce the generalizability and transferability of the approaches to other data. 

Out of the four other approaches to combine sub-quadratic runtime complexity with low data specificity, ours performs best with regards to recall and $F_{0.5}$ score, and is second in precision by a very close margin. Against this background, our achieved detection performance is comparably outstanding, given the specialized requirements in terms of scalability as well as the focus on a high-precision identification.
